% part: intuitionistic-logic
% chapter: introduction
% section: axiomatic-derivations

\documentclass[../../../include/open-logic-section]{subfiles}

\begin{document}

\olfileid{int}{int}{axd}

\olsection{مسلّمی \usetoken{P}{derivation}}

وجدانیت پسند قضیاتی منطق لئی مسلّمی !!{derivation}s خیال دے لحاظ
نال سبھ توں سادے تے تاریخ وچ سبھ توں پہلے !!{derivation} نظام نیں۔
ایہ کلاسیکی قضیاتی منطق وانگ ای کم کردے نیں۔

\begin{defn}[!!^{derivability}]
جے $\Gamma$، $\Lang L$ دے !!{formula}s دا سیٹ اے، تاں $\Gamma$
توں اک \emph{!!{derivation}}، اک محدود لڑی $!A_1$،
\dots،~$!A_n$ اے جیہڑی !!{formula}s دی اے۔ ایس وچ ہر $i \le n$ لئی ہِیٹھلیاں وچوں اک
گل لاگو اے:
\begin{enumerate}
\item $!A_i \in \Gamma$۔
\item $!A_i$ اک مسلّمہ اے۔
\item $!A_i$ کسے $!A_j$ تے $!A_k$ توں، جتھے $j < i$ تے $k <
  i$، قاعدۂ اثباتِ مقدم نال نکلدا اے، یعنی $!A_k \ident !A_j \lif !A_i$۔
\end{enumerate}
\end{defn}

\begin{defn}[مسلّمے]
وجدانیت پسند قضیاتی منطق دے \emph{مسلّمیاں} دے سیٹ $\PAx$ وچ ہِیٹھلیاں
شکلاں والے سارے !!{formula}s شامل نیں:
\begin{align}
  & (!A \land !B) \lif !A \ollabel{ax:land1}\\
  & (!A \land !B) \lif !B \ollabel{ax:land2}\\
  & !A \lif (!B \lif (!A \land !B)) \ollabel{ax:land3}\\
  & !A \lif (!A \lor !B) \ollabel{ax:lor1}\\
  & !A \lif (!B \lor !A) \ollabel{ax:lor2}\\
  & (!A \lif !C) \lif ((!B \lif !C) \lif ((!A \lor !B) \lif !C)) \ollabel{ax:lor3}\\
  & !A \lif (!B \lif !A) \ollabel{ax:lif1}\\
  & (!A \lif (!B \lif !C)) \lif ((!A \lif !B) \lif (!A \lif !C)) \ollabel{ax:lif2}\\
  & \lfalse \lif !A \ollabel{ax:lfalse1}
\end{align}
\end{defn}

\begin{defn}[!!^{derivability}]
!!^a{formula}~$!A$، $\Gamma$ توں \emph{!!{derivable}} اے،
$\Gamma \Proves !A$ لکھدے آں، جے !!a{derivation} ہووے جیہڑی
$\Gamma$ توں $!A$ اُتے مکدی اے۔
\end{defn}

\begin{defn}[قضیے]
!!^a{formula}~$!A$ اک \emph{قضیہ} اے جے !!a{derivation}~$!A$ دی
خالی سیٹ توں ہووے۔ $\Proves !A$ لکھدے آں جے $!A$ قضیہ اے،
تے $\Proves/ !A$ لکھدے آں جے نہیں۔
\end{defn}

\begin{prop}
  جے وجدانیت پسند منطق دے مسلّمی !!{derivation} نظام وچ
  $\Gamma \Proves !A$، تاں کلاسیکی منطق وچ وی $\Gamma \Proves !A$۔
  خاص کر کے، جے $!A$ وجدانیت پسند قضیہ اے، تاں کلاسیکی قضیہ وی اے۔
\end{prop}

\begin{proof}
  ہر وجدانیت پسند مسلّمہ کلاسیکی مسلّمہ وی اے۔ سو وجدانیت پسند منطق
  دی ہر !!{derivation} کلاسیکی منطق وچ وی !!a{derivation} اے۔
\end{proof}

\end{document}
